\section{Design}
\label{sec:design}
\subsection{Policy-Aligned Compact Features}
\label{sec:design-objective}
The visual interface has a teacher path and a compressed student path
during training. The teacher uses the original visual encoder and frozen
policy on clean images to produce reference tokens and actions. The
student path replaces visual processing with a compact tower, feature
encoding, quantization, and feature decoding. Its recovered tokens enter
the same policy interface. The policy weights remain fixed, but gradients
propagate through the policy to update the trainable visual path.

The objective combines three terms. The action term penalizes changes in
policy behavior after compression. For $\pi_{0.5}$, training uses the
frozen policy's flow-matching objective with targets derived from its own
clean-input actions. For OpenVLA-OFT, the recorded action-distillation
configuration uses an L1 loss on clean-policy action chunks. The token
term aligns recovered tokens with the original visual representation
using mean-squared error and cosine similarity. The rate term penalizes
coding cost and is adjusted toward a target payload during
training. These terms respectively constrain the downstream behavior,
the intermediate representation, and the communication cost.

The clean actions in this procedure are outputs of the original policy,
not additional expert action annotations. Token supervision anchors the
student to the original interface, while action supervision penalizes
representation errors through their effect on the frozen policy. The two
terms are complementary design constraints, rather than two independent
deployment heads. Only the compressed path is needed at inference time.
The rate penalty is a training objective; actual integer payloads are
measured separately after entropy coding.

\paragraph{Compact edge representation.}
The primary $\pi_{0.5}$ path distills a SigLIP visual
teacher~\cite{siglip}. In our implementation, the student has
12 transformer layers with width 768 and approximately 87 million
parameters, compared with 415 million parameters for the selected
teacher configuration. A transformer feature encoder uses a
256-dimensional bottleneck. The main configuration uses a
five-level scalar quantizer with step size 4 and bound 8. All camera
streams pass through the learned feature path; the main configuration
does not substitute a JPEG stream for one of the cameras.

\paragraph{Direct recovery of policy tokens.}
A four-layer feature decoder with width 512 recovers 256 tokens of
dimension 2,048 per camera in the $\pi_{0.5}$ configuration. These tokens
replace the outputs of the original visual tower and its projection into
the policy. Image reconstruction is absent from the deployed cloud path.
The decoder is trained for the interface expected by the frozen policy,
so its output is not intended as a general-purpose visual representation
for unrelated consumers.

\subsection{Decoder-Free Temporal Coding at the Edge}
\label{sec:design-temporal}
Successive uploads can contain repeated visual information. We exploit
this redundancy with differential pulse-code modulation (DPCM) in the
quantized feature representation. The prediction reference is the
previous quantized code, which is retained at both endpoints. An intra
upload establishes a reference; inter uploads encode a residual with
respect to that reference. The primary $\pi_{0.5}$ configuration uses a
group of pictures (GOP) of 16 uploads, and the OFT configuration uses a
GOP of 8. These intervals refer to policy uploads rather than every
camera capture or actuator step.

The reference choice determines the local work. If prediction depended
on the output of a neural reconstruction decoder, the edge would need
to execute that decoder to reproduce the cloud's reference. Using the
quantized code instead lets the sender maintain its reference without a
neural decoding pass. The receiver first reverses entropy and temporal
coding, then uses the recovered code as input to the feature decoder.
Both endpoints retain the corresponding quantized reference for the
next upload.

This arrangement separates discrete stream correctness from model
numerics. The entropy coder and temporal reconstruction must reproduce
the transmitted quantized representation exactly; the recovered policy
tokens are still a learned approximation to the original visual
features. Likewise, identical entropy coding across ARM and x86 hosts
does not imply that floating-point encoders on different devices always
produce identical symbols for the same image.

\paragraph{Integer bitstream.}
We implement static rANS~\cite{ans} with a 32-bit state and integer
frequency tables normalized to 15 bits. The decoder uses the matching
tables to recover symbols. This supplies a concrete byte stream rather
than relying only on the rate estimate used for training. Payload
accounting must include the components actually needed by the selected
deployment; application framing and network traffic are distinct from
the codec payload.

\paragraph{Temporal state.}
Training rolls through intra and inter uploads to expose the codec to
the reference states used by deployment. Streams must begin from a
consistent reference and advance in the same order at both endpoints.
Periodic intra coding limits the interval between reference refreshes,
but it does not by itself establish packet-loss or reconnection
robustness. Our measurements do not validate a general recovery protocol
for corrupted or missing reference state.

\subsection{Integration with Frozen VLA Backbones}
\label{sec:design-integration}
We adapt the feature interface to two backbones rather than imposing a
single token format on both. For $\pi_{0.5}$, the recovered representation
replaces SigLIP and its projection. OpenVLA uses visual features derived
from DINOv2 and SigLIP~\cite{openvla,dinov2,siglip}; the OFT student and
decoder are adapted to this dual-tower interface. The OFT student
has width 1,024. Table~\ref{tab:interfaces} summarizes the scope of these
adaptations.

\begin{table}[t]
\caption{Backbone-specific visual interfaces. Each codec is trained for
its associated policy; policy weights remain frozen.}
\label{tab:interfaces}
\centering\small
\begin{tabularx}{\columnwidth}{lXX}
\toprule
Backbone & Teacher visual path & Recovered representation \\
\midrule
$\pi_{0.5}$ & SigLIP and projection & Policy visual tokens \\
OpenVLA-OFT & DINOv2 and SigLIP & Adapted dual-tower features \\
\bottomrule
\end{tabularx}
\end{table}

The shared components are student-based encoding, a quantized temporal
stream, and a cloud feature adapter. Their shapes and supervision are
policy dependent. The action-distillation objective also follows the
policy's action parameterization. This makes the approach adaptable to
multiple interfaces, but not a zero-shot codec transferable between
policies without retraining.
